% from ega4-18-fr.tex
\oldpage[IV-4]{124}

\begin{corollary}[18.4.12]
\label{IV.18.4.12}
\begin{enumerate}
  \item[(i)] Soient $f:X\to Y$ un morphisme localement de type fini, $x$ un point de $X$. Si $f$ est plat et formellement non ramifié au point $x$, alors, dans toute factorisation $f|U:U\xrightarrow{j}X'\xrightarrow{h}Y$, où $U$ est un voisinage ouvert de $x$, $j$ une immersion fermée et $h$ un morphisme étale (18.4.7), l'homomorphisme $\sh{O}_{X',j(x)}\to\sh{O}_{X,x}$ correspondant à $j$ est bijectif (en particulier $\sh{O}_{X,x}$ est une $\sh{O}_{Y,f(x)}$-algèbre \emph{essentiellement étale} (18.6.1)).
  \item[(ii)] Pour qu'un morphisme $f:X\to Y$ soit étale, il faut et il suffit qu'il soit localement de type fini, formellement non ramifié et plat.
\end{enumerate}
\end{corollary}

(i) Comme l'homomorphisme $\sh{O}_{X',j(x)}\to\sh{O}_{X,x}$ est surjectif, il suffit de prouver qu'il est injectif, et pour cela qu'il fait de $\sh{O}_{X,x}$ un $\sh{O}_{X',j(x)}$-module \emph{fidèlement plat}, ou seulement \emph{plat} ($0_{\mathrm I}$, 6.5.1 et 6.6.2); autrement dit, il s'agit de montrer que $j$ est un morphisme \emph{plat au point} $x$; mais puisque $h\circ j=f$ est par hypothèse plat au point $x$ et que $h$ est étale, cela résulte de (18.4.10) et (18.4.11, (i)).

(ii) Il n'y a à prouver que la suffisance des conditions énoncées. Pour tout $x\in X$, on a donc dans un voisinage ouvert $U$ de $x$ une factorisation de $f|U$ ayant les propriétés considérées dans (i). Comme par hypothèse $f$ est plat et formellement non ramifié en \emph{tous} les points de $U$, le résultat de (i) s'applique non seulement à $x$ mais à tous les points de $U$; cela signifie que si $\mathscr J$ est l'Idéal quasi-cohérent de $\sh{O}_{X'}$ correspondant au sous-préschéma fermé de $X'$ associé à $j$, on a $\mathscr J_{j(z)}=0$ pour tout $z\in U$, donc $j$ est une immersion ouverte, et $f$ est par suite étale en tout point de $U$, donc en tout point de $X$.

\begin{corollary}[18.4.13]
\label{IV.18.4.13}
Soient $f:X\to Y$ un morphisme localement de type fini, $x$ un point de $X$, $y=f(x)$. Supposons que $y$ admette un voisinage ouvert qui soit un préschéma réduit et n'ayant qu'un nombre fini de composantes irréductibles. Alors, pour que $f$ soit étale au point $x$, il faut et il suffit que $f$ soit plat et formellement non ramifié au point $x$.
\end{corollary}

Il n'y a à démontrer que la suffisance de la condition. La question étant locale sur $X$ et $Y$, on peut supposer (18.4.7) que $f$ se factorise en $X\xrightarrow{j}X'\xrightarrow{h}Y$ où $h$ est étale et $j$ une immersion fermée, et en outre que $Y$ est réduit et n'a qu'un nombre fini de composantes irréductibles. Alors $X'$ est réduit (17.5.7) et, en remplaçant au besoin $X'$ par un voisinage ouvert de $j(x)$, on peut supposer que $X'$ n'a qu'un nombre fini de composantes irréductibles: en effet, on peut supposer $h$ quasi-fini (17.6.1) et comme les points maximaux de $X'$ sont au-dessus des points maximaux de $Y$ (2.3.4) leur nombre est fini. Tout revient à montrer, avec les notations de la démonstration de (18.4.12), que l'on a $\mathscr J_{x'}=0$ pour tous les points $x'$ d'un voisinage de $j(x)$ dans $X'$, sachant que $\mathscr J_{j(x)}=0$. Or, en remplaçant $X'$ par un voisinage affine de $j(x)$, on peut supposer que toutes les composantes irréductibles de $X'$ contiennent $j(x)$; si $X'=\Spec(A')$, et si $\mathfrak p'$ est l'idéal premier de $A'$ correspondant au point $j(x)$, le morphisme $\Spec(A'_{\mathfrak p'})\to\Spec(A')$ est dominant, donc l'homomorphisme correspondant $A'\to A'_{\mathfrak p'}$ est \emph{injectif} puisque $A'$ est réduit (I, 1.2.7). Si $\mathscr J=\widetilde{\mathfrak J}$, où $\mathfrak J$ est un idéal de $A'$, $\mathfrak J$ s'identifie donc à une partie de $\mathfrak J_{\mathfrak p'}$, et l'hypothèse $\mathfrak J_{\mathfrak p'}=0$ entraîne donc $\mathfrak J=0$.

\begin{corollary}[18.4.14]
\label{IV.18.4.14}
Soient $A$ un anneau local d'idéal maximal $\mathfrak m$, de corps résiduel $k$, $B$ une $A$-algèbre finie.

